\section{Bối cảnh thực trạng}

Với sự phát triển nhanh chóng của kinh tế độ cao thấp, giao hàng bằng drone đã nổi lên như một giải pháp logistics đầy hứa hẹn nhờ tốc độ, tính linh hoạt và khả năng tiếp cận các khu vực xa xôi. Theo ResearchAndMarkets.com, thị trường giao hàng bằng drone toàn cầu được định giá khoảng 5 tỷ USD vào năm 2024 và dự kiến sẽ đạt 33,4 tỷ USD vào năm 2030, với tốc độ tăng trưởng kép hằng năm (CAGR) 37,3\%. Sự tăng trưởng này phản ánh nhu cầu ngày càng cao về các giải pháp giao hàng chặng cuối nhanh chóng, hiệu quả và thân thiện với môi trường trong lĩnh vực thương mại điện tử và y tế.

Tuy nhiên, ngoài những thách thức về phần cứng và quy định, trở ngại lớn nhất nằm ở việc quản lý và mở rộng hoạt động giao hàng. Việc phối hợp nhiều trạm hạ cánh, theo dõi tình trạng đơn hàng theo thời gian thực, và ước tính chính xác thời gian giao hàng trên một đội drone phân tán đòi hỏi công cụ vận hành chuyên biệt vượt xa khả năng điều khiển bay đơn thuần. Các hệ thống quản lý logistics hiện có thường được thiết kế cho giao hàng mặt đất và thiếu hỗ trợ cho các quy trình đặc thù của drone như kiểm tra tình trạng trạm hạ cánh, xác nhận lấy hàng, hay ước tính thời gian giao hàng bằng AI.

Do đó, các doanh nghiệp áp dụng giao hàng bằng drone cần một nền tảng quản lý chuyên biệt có khả năng tự động hóa các quy trình này và tối ưu hóa toàn bộ vòng đời giao hàng -- từ khâu tạo đơn đến xác nhận giao hàng.

